\documentclass[11pt]{article}

\usepackage[margin=1in]{geometry}
\usepackage{amsmath,amssymb}
\usepackage{hyperref}
\usepackage{enumitem}
\usepackage{xcolor}

\hypersetup{
  colorlinks=true,
  linkcolor=blue,
  urlcolor=blue,
  citecolor=blue
}

\title{Proto-Teleology Without Purpose:\\
Consequence-Bearing Difference and Sound Identification}
\author{S. Poole}
\date{June 12, 2026}

\newcommand{\Rel}{\mathrm{Rel}}
\newcommand{\Sep}{\mathrm{Sep}}
\newcommand{\Asym}{\mathrm{Asym}}

\begin{document}

\begin{titlepage}
\centering
\vspace*{1.2in}
{\Huge\bfseries Proto-Teleology Without Purpose\par}
\vspace{0.35in}
{\LARGE Consequence-Bearing Difference and Sound Identification\par}
\vspace{0.85in}
{\Large Citation-Oriented Internal Review Draft\par}
\vfill
{\large S. Poole\par}
\vspace{0.2in}
{\large June 12, 2026\par}
\vspace{0.7in}
{\small Draft manuscript for internal review.\par}
\end{titlepage}

\begin{abstract}
This paper introduces a minimal formal account of proto-teleology: directed
consequence without purpose, value, agency, selfhood, or moral truth. The
guiding question is simple: when does a difference matter because erasing it
changes what can follow? We answer in two stages. First, a primitive relational
presentation supplies a non-collapse condition: some relata are connected,
distinguished, and asymmetrically arranged in a way that cannot be flattened
into total interchangeability. Second, an admissible continuation apparatus
tests whether that primitive difference changes consequences under declared
contexts. When an evaluated context refuses the erasure of a primitive
difference, the difference becomes consequence-bearing. We call this minimal
structure proto-teleological because it introduces directional non-neutrality
before any notion of goal, utility, organism, agent, or valuer.

The framework treats labels, quotient classes, clusters, and learned
representations as downstream presentations. They are acceptable only when
their proposed identifications respect consequence-induced separation. This
aligns the proposal with standard concerns in quotienting, contextual
equivalence, abstraction soundness, and bisimulation-like reasoning, while
remaining lower than any claim about value or agency. We also describe finite
baseline witnesses: small worlds where familiar summaries agree while declared
consequence or recovery facts differ. These witnesses show why exact guardrails
are needed before scaling to richer substrates.
\end{abstract}

\section{Introduction}

Many theories of value, agency, and purpose begin with something already
structured: a subject, a utility function, a reward signal, an organism, a
self-maintaining system, a boundary, a preference ordering, or a goal. That is
often appropriate for applied modeling, and it has deep roots in cybernetic and
biological discussions of purpose and teleology
\cite{rosenblueth1943behavior,mayr1961cause,monod1971chance}. But it leaves a
prior question open:

\begin{quote}
What makes any difference matter before value or valuers are assumed?
\end{quote}

This paper addresses that prior question. It does not try to derive morality,
agency, or value. It tries to formalize the first anti-flattening step: a
difference matters in the minimal technical sense when erasing it changes
admissible continuation.

The central claim is:

\begin{quote}
A primitive difference becomes proto-teleological when admissible continuation
exposes its erasure as consequence-changing.
\end{quote}

Here ``proto-teleological'' means directed consequence before purpose. It is a
weak and deliberately constrained term. It says that some continuations,
identifications, or erasures are not neutral. It does not say that the system
wants anything.

\subsection{Why this matters}

This lower layer matters for four reasons.

First, it gives a disciplined sense of ``mattering'' before value. A difference
can be consequence-bearing even in a substrate with no preferences, no agents,
and no moral facts. That is not enough for ethics, but it blocks a premature
flattening move where all distinctions are treated as arbitrary unless a valuer
already exists.

Second, it disciplines abstraction. In science and machine learning, it is
common to compress states, cluster observations, learn representations, or
quotient a system by apparent similarity. The formal lesson here is that a
compression is not trustworthy merely because it is concise or predictive. It
must be checked against the consequences it merges.

Third, it clarifies a possible route from non-neutrality to richer theories of
value-bearing systems. If consequence-bearing differences can recur, compose,
survive abstraction, support compatible continuations, and eventually support
valuer-like trajectories, then value need not be injected at the primitive
layer. This paper does not prove that route. It identifies one lower condition
such a route would have to preserve.

Fourth, it makes small finite worlds useful. In finite settings, profiles can
be exhaustive, failures can be exact, and negative controls can be checked
without appeal to intuition. That gives us a way to test reductions before
making claims about large or physical substrates.

\subsection{Contributions}

\begin{enumerate}[leftmargin=*]
  \item We isolate a primitive non-collapse condition using relation,
  distinction, and asymmetry.

  \item We define a continuation/consequence apparatus as an operational test
  of whether erasure is neutral.

  \item We define proto-teleology as a consequence filter: evaluated contexts
  may allow or refuse proposed identifications.

  \item We connect the framework to established formal ideas: quotient
  soundness, contextual equivalence, abstraction contracts, and
  bisimulation-style guardrails.

  \item We describe finite baseline witnesses showing that familiar summaries
  can agree while declared consequence or recovery facts differ.
\end{enumerate}

\section{Terminology in Plain Language}

This section fixes the working vocabulary. The goal is not to introduce a new
dictionary but to make precise which standard idea each term is trying to
track.

\subsection{Difference, distinction, and erasure}

A \emph{difference} is an informal contrast between two possible fragments or
relata. A \emph{distinction} is a formal predicate witnessing that two relata
are not freely interchangeable in the primitive presentation.

An \emph{erasure} is the act of identifying two relata, treating them as the
same, or placing them in the same quotient class. In formal implementations
this often appears as a proposed merge.

The central question is:

\begin{quote}
Is this erasure allowed by the declared continuation consequences?
\end{quote}

\subsection{Continuation and consequence}

A \emph{continuation} is any admissible way the substrate can be extended,
evolved, completed, transformed, derived, composed, perturbed, or probed. Time
evolution is only one kind of continuation.

A \emph{consequence} is what a fragment yields under a continuation context. In
a transition system this may be a reachable future. In a graph it may be
reachability. In a logic it may be derivability. In a stochastic channel it may
be an output distribution. In a constraint system it may be completion
behavior.

\subsection{Compatibility, separation, and identification}

Two fragments are \emph{compatible} under a declared panel if every evaluated
context compares their consequences successfully. A context \emph{separates}
them when it refuses that comparison.

Symmetric \emph{identification} is stronger than one-way compatibility. If the
comparison relation is directional, then allowing \(a\) as a substitute for
\(b\) does not automatically allow \(b\) as a substitute for \(a\).

This is close in spirit to contextual equivalence and observational
equivalence: two things may be identified only relative to the contexts that
fail to distinguish them. The difference here is that we keep the comparison
relation and evaluated panel explicit, because a bad panel can create false
collapse or false separation.

\subsection{Filter and gradient}

The current formal layer establishes a \emph{filter}, not a full gradient.

A filter says:

\begin{quote}
this erasure is allowed; that erasure is refused.
\end{quote}

A gradient would additionally order continuations or profiles:

\begin{quote}
this continuation preserves more declared consequence-bearing structure; that
one erases more; this profile dominates that profile.
\end{quote}

The filter is currently formalized. The gradient remains future work.

\section{Related Work and Positioning}

The proposal sits between several established literatures, but it is not a
replacement for any of them.

The term proto-teleology is chosen against the background of cybernetic and
biological accounts of purpose. Rosenblueth, Wiener, and Bigelow analyzed
purposeful behavior in terms of feedback and observable behavior
\cite{rosenblueth1943behavior}; Mayr distinguished different kinds of causal
explanation in biology \cite{mayr1961cause}; and Monod popularized
``teleonomy'' as a way to discuss apparent purposiveness in living systems
without appealing to final causes \cite{monod1971chance}. The present paper is
lower-level than those accounts. It does not require an organism, a feedback
controller, or an evolved biological function. It asks when a distinction is
already non-neutral under admissible continuation.

The quotient and identification discipline is closer to formal methods.
Bisimulation and coinduction study when transition systems can be treated as
behaviorally equivalent \cite{milner1989communication,sangiorgi2012bisimulation}.
Probabilistic bisimulation and Markov-chain lumpability similarly constrain
which stochastic states can be identified
\cite{larsen1991bisimulation,kemeny1960finite}. Model checking emphasizes
finite-state structures and specification-relative verification
\cite{baier2008principles}. The present framework borrows the spirit of these
methods: identifications are justified by preservation of declared behavior,
not by labels. Its contribution is to make the context panel and comparison
relation explicit enough to talk about consequence-bearing difference before
committing to one particular transition-system formalism.

The abstraction discipline is also related to abstract interpretation, where an
abstract semantics is useful only when it soundly approximates a concrete one
\cite{cousot1977abstract}. This paper uses the same broad moral: coarse views
must be checked against exact facts. It also intersects information-theoretic
compression, especially Shannon's founding account of information
\cite{shannon1948mathematical} and the information bottleneck method's
compression/prediction tradeoff \cite{tishby1999information}. The difference is
that the present target is not compression alone. A summary can be compact or
predictive while still merging fragments that a declared consequence apparatus
separates.

\section{The Primitive Floor: Non-Collapse}

The primitive presentation has three roles:

\begin{description}[leftmargin=*]
  \item[Relation] Relata stand in structured bearing.
  \item[Distinction] Two relata are not freely interchangeable under some
  formal contrast.
  \item[Asymmetry] The bearing of a distinction is not neutral in all
  directions.
\end{description}

Formally, these can be represented by predicates analogous to:

\[
\Rel(x,y), \qquad \Sep(d,x,y), \qquad \Asym(d,x,y).
\]

The theory does not treat those field names as magic. Empty predicates,
decorative predicates, or disconnected predicates do not provide a real
primitive floor. The stack therefore distinguishes vocabulary from witnesses.

\subsection{Joint witnesses}

A joint primitive witness says relation and separation touch the same relata:

\[
\exists d,x,y.\ \Rel(x,y) \wedge \Sep(d,x,y).
\]

This blocks two total collapses:

\begin{description}[leftmargin=*]
  \item[RelationCollapse] No relation holds anywhere.
  \item[IdentificationCollapse] No distinction separates any two relata.
\end{description}

\subsection{Asymmetry witnesses}

An asymmetry witness says:

\[
\exists d,x,y.\ \Asym(d,x,y).
\]

In the formal stack used by this paper, asymmetry implies both relation and
separation. So an asymmetry witness supplies a joint witness. This factoring
matters: the collapse-blocking work is done by relation/separation contact;
asymmetry is a strong source of that contact.

The result is narrow but important:

\begin{quote}
An actual primitive witness prevents total flattening.
\end{quote}

It does not yet say that the distinction changes what can follow. That requires
continuation consequence.

\section{Continuation as a Contextual Test}

The continuation layer asks whether a primitive difference is operationally
exposed. It uses a consequence system with:

\begin{description}[leftmargin=*]
  \item[Fragment] The item being compared.
  \item[Context] An admissible continuation or test condition.
  \item[Outcome] What the fragment yields under that context.
  \item[consequence] A map \(\mathrm{Context} \to \mathrm{Fragment} \to
  \mathrm{Outcome}\).
  \item[Compare] Context-indexed comparison of outcomes.
  \item[Evaluated] The contexts actually in the declared panel.
\end{description}

This structure is intentionally close to familiar ideas:

\begin{itemize}[leftmargin=*]
  \item Contextual equivalence: two terms are equivalent if no context
  distinguishes them.

  \item Observational equivalence: two states are equivalent if observations
  cannot tell them apart.

  \item Bisimulation and quotienting: state merges are valid only when the
  transition or observation structure is respected.

  \item Abstraction soundness: a coarse model is acceptable only if its claims
  remain true relative to the exact system.
\end{itemize}

The proposal does not simply import one of these frameworks because it wants a
lower and more general interface. The consequence system can represent temporal
dynamics, graph reachability, logical derivability, stochastic output,
constraint completion, or other continuation forms.

\subsection{Compatibility and separation}

Directional compatibility means:

\[
\mathrm{Compatible}(a,b) \equiv
\forall k.\ \mathrm{Evaluated}(k) \to
\mathrm{Compare}_k(\mathrm{consequence}(k,a),\mathrm{consequence}(k,b)).
\]

Directional separation means:

\[
\mathrm{Separated}(a,b) \equiv
\exists k.\ \mathrm{Evaluated}(k) \wedge
\neg \mathrm{Compare}_k(\mathrm{consequence}(k,a),\mathrm{consequence}(k,b)).
\]

If comparison is directional, separation is directional. Therefore the formal
stack distinguishes:

\begin{description}[leftmargin=*]
  \item[one-way allowance] \(a\) can stand in for \(b\) under the declared
  comparisons;
  \item[symmetric identification] \(a\) can stand in for \(b\) and \(b\) can
  stand in for \(a\);
  \item[merge separation] at least one direction is separated, so symmetric
  identification is blocked.
\end{description}

The core rule is:

\begin{quote}
If evaluated consequence merge-separates \(a\) and \(b\), then \(a\) and \(b\)
cannot be soundly identified.
\end{quote}

\section{Proto-Teleology as a Consequence Filter}

A proto-teleological condition exists when:

\begin{enumerate}[leftmargin=*]
  \item there is a primitive witness, and
  \item evaluated consequence merge-separates that witness's endpoints.
\end{enumerate}

Equivalently:

\[
\text{primitive non-collapse}
+
\text{consequence refusal}
=
\text{consequence-bearing difference}.
\]

This is the first formal sense in which a difference matters. The difference
does not matter because a subject values it. It matters because erasing it
changes what the declared continuation apparatus permits.

\subsection{The exact profile}

The consequence profile records which pairs are blocked and which are allowed:

\[
\mathrm{Blocks}(a,b) \equiv
\text{evaluated consequence merge-separates } a \text{ and } b,
\]

\[
\mathrm{Allows}(a,b) \equiv
a \text{ and } b \text{ are symmetrically identifiable}.
\]

The checked bridge is:

\[
\mathrm{ProtoTeleologicalCondition}(S)
\to
\exists a,b.\ \mathrm{Blocks}_S(a,b).
\]

So the proto-teleological condition induces a nonempty blocked-merge profile.
That is the filter.

\subsection{Why this is not value}

This result does not rank outcomes. It does not say a continuation is better.
It does not say the blocked merge is morally wrong. It says only that a
declared consequence test refuses an erasure.

The philosophical significance is lower:

\begin{quote}
Non-neutrality precedes valuation.
\end{quote}

If no differences can be consequence-bearing, value has nowhere to attach. But
if consequence-bearing differences exist, the later question of value becomes
well-posed without being answered.

\section{Why Quotients and Labels Are Downstream}

A recurring failure mode in formal and empirical work is to start with a
classification and then treat the classification as ontological. The present
stack reverses that order.

Do not start from:

\begin{itemize}[leftmargin=*]
  \item cluster labels;
  \item learned representation coordinates;
  \item quotient classes;
  \item visual similarity;
  \item predictive success alone.
\end{itemize}

Start from:

\begin{quote}
which proposed identifications are refused by evaluated consequence?
\end{quote}

This is a quotient-soundness requirement. A quotient is valid only if every
merge it performs is allowed by the consequence profile. If a quotient class
contains a separated pair, it is invalid relative to the declared apparatus.

\subsection{Chain-connectedness is not identity}

One common mistake is to infer class membership from chains:

\[
a \sim b,\quad b \sim c,\quad \therefore a \sim c.
\]

That inference is valid only under additional structure, such as transitivity.
If the comparison relation is nontransitive, \(a\) may compare with \(b\), and
\(b\) may compare with \(c\), while \(a\) is separated from \(c\).

For nontransitive comparisons, a proposed class must be pairwise compatible
unless stronger structure has been proved. This is a small theorem with a large
methodological consequence: connectedness is not identity.

\section{Abstraction and Coarse-Graining}

Coarse-graining is necessary in large systems. The point is not to reject it.
The point is to make it accountable.

The current abstraction layer separates exact profiles from coarse claims.
Given an exact consequence system, an abstraction may claim:

\[
\mathrm{allows}(a,b), \qquad \mathrm{blocks}(a,b).
\]

Those claims are judged by two independent contracts:

\begin{description}[leftmargin=*]
  \item[soundness] abstraction claims imply exact profile facts;
  \item[completeness] exact profile facts imply abstraction claims.
\end{description}

An abstraction can be sound but incomplete. It can also be complete but
unsound. This distinction matters because a coarse summary can look clean while
erasing the exact consequence facts that made the structure relevant.

This is the connection to scale. We should not claim that all coarse-grainings
preserve primitive or consequence structure. We should ask which
transformations and abstractions preserve the facts required by the theory.

\section{Finite Baseline Witnesses}

Finite worlds let us make exact mistakes on purpose. That is useful.

The baseline witness suite contains finite examples where a familiar summary
matches across two systems while a declared consequence or recovery fact
differs. The retained patterns include:

\begin{itemize}[leftmargin=*]
  \item same reachability, different declared recovery;
  \item same mutual information, different declared recovery;
  \item same chain evidence, different class soundness;
  \item same compression score, different merge soundness;
  \item same coarse bisimulation, different consequence profile.
\end{itemize}

The suite supports this narrow claim:

\begin{quote}
Familiar finite summaries can match while declared consequence/recovery,
merge-soundness, or class-soundness facts differ.
\end{quote}

This matters because it attacks a specific reduction strategy. If someone says
``your structure is just mutual information,'' or ``just reachability,'' or
``just bisimulation,'' the finite witnesses show that those summaries can agree
while the declared consequence facts disagree.

The finite witnesses do not prove that all reductions fail. They do not prove
physical or biological relevance. They provide exact warning signs and a
reproducible adversarial test surface.

\section{Validation Discipline}

The current validation stack supports three kinds of checks.

\subsection{Formal proofs}

The Lean stack checks primitive nondegeneracy, collapse blockers, consequence
separation, class guardrails, panel discipline, proto-teleological conditions,
profile bridges, abstraction contracts, and selected finite baseline
witnesses.

\subsection{Finite search and smoke tests}

Python runners check retained finite witness examples and parameterized finite
families. These are not substitutes for theorems, but they make external
review and adversarial search easier.

\subsection{CI and reproducibility}

Cross-platform smoke checks make the finite witness suite easier to validate
outside the development machine. The larger formal build checks the active
proof umbrella.

The validation principle is:

\begin{quote}
Every positive structure should be paired with failure modes.
\end{quote}

Vacuous contexts, universal comparisons, all-refusing panels, unsound
abstractions, and reduction baselines are not side issues. They are the
discipline that prevents the definitions from becoming self-validating.

\section{Larger Ambition}

The present paper establishes only the lower floor:

\[
\text{primitive non-collapse}
\to
\text{evaluated consequence refusal}
\to
\text{consequence-bearing difference}
\to
\text{exact allow/block filter}.
\]

A fuller theory of value-bearing continuation would require much more:

\[
\text{filter/profile}
\to
\text{gradient or preorder over profiles}
\to
\text{consequence-preserving transformations}
\to
\text{recurrence}
\to
\text{recoverability without exact identity}
\to
\text{support or extent}
\to
\text{compatibility over longer continuations}
\to
\text{valuer-capable trajectories}.
\]

The reason this lower floor matters is that it gives that larger program a
non-arbitrary starting point. Instead of assuming value, the proposed path asks
how consequence-bearing differences can recur, compose, survive abstraction,
and eventually support valuers.

That larger program remains open. The paper earns only the first step.

\section{Limitations}

\begin{enumerate}[leftmargin=*]
  \item Primitive non-collapse alone does not prove proto-teleology.
  Proto-teleology requires consequence responsiveness under an admissible
  continuation apparatus.

  \item Apparatus admissibility remains central. Arbitrary contexts and
  comparisons can manufacture separations or collapses.

  \item The current structure is a filter, not a gradient. Ordering
  continuations or profiles remains future work.

  \item Recoverability is not established. Exact identity is too strong; future
  work needs transformation and abstraction contracts.

  \item Finite baseline witnesses are not infinite-family theorems and do not
  prove substrate-general validity.

  \item No claim is made about value, morality, agency, selfhood, life,
  identity, boundary, deformer structure, or any terminal theory of
  value-bearing futures.
\end{enumerate}

\section{Conclusion}

The paper's core claim is modest but useful:

\begin{quote}
A difference matters, in the proto-teleological sense, when erasing it changes
admissible continuation.
\end{quote}

This gives a formal foothold for directed consequence without purpose. It does
not smuggle in a valuer, a goal, a self, or a utility function. It says only
that some primitive differences become consequence-bearing under declared
continuation tests.

That is enough to block total flattening. It is not enough to validate a theory
of value. It is the first layer from which such a theory can be built without
making value primitive.

\appendix

\section{Repository Pointers}

The repository paths below use project-internal implementation names. They are
not part of the paper's conceptual vocabulary.

Public onboarding documents:

\begin{itemize}[leftmargin=*]
  \item \texttt{docs/OMEGA\_FORMALISM\_PRIMER.md}
  \item \texttt{docs/EXTERNAL\_READER\_GUIDE.md}
  \item \texttt{docs/CLAIMS\_LEDGER.md}
  \item \texttt{docs/BASELINE\_WITNESS\_SUITE\_V0.md}
\end{itemize}

Relevant Lean files:

\begin{itemize}[leftmargin=*]
  \item \texttt{formal/lean/AlphaCore/Primitive.lean}
  \item \texttt{formal/lean/AlphaCore/Nondegenerate.lean}
  \item \texttt{formal/lean/AlphaCore/PrimitiveMap.lean}
  \item \texttt{formal/lean/OmegaProper/Trajectory/ConsequenceRelation.lean}
  \item \texttt{formal/lean/OmegaProper/Trajectory/ConsequenceClasses.lean}
  \item \texttt{formal/lean/OmegaProper/Trajectory/ConsequenceDiscipline.lean}
  \item \texttt{formal/lean/OmegaProper/Trajectory/ConsequenceComparison.lean}
  \item \texttt{formal/lean/OmegaProper/Trajectory/AlphaConsequenceSeed.lean}
  \item \texttt{formal/lean/OmegaProper/Trajectory/ProtoTeleologicalSeed.lean}
  \item \texttt{formal/lean/OmegaProper/Trajectory/ProtoTeleologicalProfile.lean}
  \item \texttt{formal/lean/OmegaProper/Trajectory/ProfileAbstraction.lean}
  \item \texttt{formal/lean/OmegaProper/BaselineWitnesses.lean}
\end{itemize}

\section{Reproduction Commands}

Lean build:

\begin{verbatim}
powershell -ExecutionPolicy Bypass -File scripts\setup\invoke_lake.ps1 build AlphaOmega
\end{verbatim}

Retained finite witness smoke:

\begin{verbatim}
python -m omega.validation.baseline_witness_smoke
\end{verbatim}

Parameterized finite witness family smoke:

\begin{verbatim}
python -m omega.validation.baseline_witness_family_smoke
\end{verbatim}

\bibliographystyle{plain}
\bibliography{proto_teleology_without_purpose_cited_draft}

\end{document}
